\section{Introduction}

\label{sec:introduction}
%% ============================================================================

Shor's algorithm~\cite{shor1994} breaks the elliptic curve discrete logarithm
problem in polynomial time on a cryptographically relevant quantum computer
(CRQC). Every BLS, ECDSA, and Schnorr-based blockchain is therefore vulnerable
to a harvest-now-decrypt-later attack: an adversary records signed blocks today
and forges certificates once a CRQC becomes available. NIST estimates this
threshold at 2030--2035~\cite{mosca2018}.

The standard response is to replace classical signatures with a post-quantum
scheme. This is insufficient for two reasons.

First, no single post-quantum signature scheme has survived decades of
cryptanalysis. ML-DSA (Dilithium)~\cite{nist-pqc} was standardized in 2024;
lattice-based threshold signing remains non-standardized. Relying on a single
post-quantum assumption is exchanging one single point of failure for another.

Second, post-quantum signatures are large. ML-DSA-65 signatures are
$3{,}309$~bytes versus 48~bytes for a BLS aggregate. Per-block ML-DSA
certificates would increase block header size by two orders of magnitude,
degrading network throughput.

Quasar resolves both problems. Three independent cryptographic layers---BLS12-381,
Ringtail (Ring-LWE threshold), and ML-DSA-65---run in parallel during each
consensus round. BLS provides fast classical finality; Ringtail provides
post-quantum threshold finality; ML-DSA-65 provides post-quantum identity
binding. The critical path is the slowest layer, not their sum, because all
three signing operations execute concurrently.

The three hardness assumptions are:
\begin{enumerate}[leftmargin=2em]
\item \textbf{Discrete logarithm} over BLS12-381 (classical, broken by Shor).
\item \textbf{Module-LWE} over $\mathbb{Z}_q[X]/(X^N+1)$ (ML-DSA-65, FIPS~204).
\item \textbf{Module-SIS / Ring-LWE} (Ringtail threshold signing).
\end{enumerate}
These assumptions are algebraically independent: no known reduction connects
the hardness of one to the hardness of another. An adversary that breaks the
discrete logarithm problem (BLS) gains no advantage against lattice problems
(Ringtail, ML-DSA), and vice versa. An adversary that breaks Module-LWE
(ML-DSA) does not automatically break Ring-LWE with different parameters
(Ringtail).

\subsection{Contributions}

\begin{enumerate}[leftmargin=2em]
\item A three-layer consensus protocol where safety holds if at least one
      cryptographic layer remains secure.
\item A concrete construction combining BLS aggregate signatures, Ring-LWE
      threshold signatures, and ML-DSA-65 identity signatures with parallel
      execution.
\item Four theorems: safety, liveness, Quasar unforgeability, and
      single-compromise insufficiency, with references to machine-checked
      proofs in Lean~4 and Tamarin.
\item Production benchmarks on commodity hardware showing $69\,\text{ms}$
      full threshold rounds at $n=100$.
\end{enumerate}


%% ============================================================================
